%%=============================================================================
%% Voorwoord
%%=============================================================================

\chapter*{\IfLanguageName{dutch}{Woord vooraf}{Preface}}%
\label{ch:voorwoord}

%% TODO:
%% Het voorwoord is het enige deel van de bachelorproef waar je vanuit je
%% eigen standpunt (``ik-vorm'') mag schrijven. Je kan hier bv. motiveren
%% waarom jij het onderwerp wil bespreken.
%% Vergeet ook niet te bedanken wie je geholpen/gesteund/... heeft

Toen ik werd aangenomen bij VLAOI, was de keuze voor dit onderwerp vrij snel gemaakt. 
Bij de eerste gesprekken kreeg ik twee concrete onderzoeksvoorstellen. Het meer data-ge\-rich\-te voorstel sprak mij daarbij aan. Het sloot aan bij mijn interesse en had een concrete meerwaarde voor de organisatie. Een proof-of-concept uitwerken dat VLAIO effectief kon gebruiken, trok mij meer aan dan een academisch experiment dat enkel voor mijzelf iets opleverde.\\

Het gaf mij zowel de kans om bij te leren over data governance, metadataplatformen en werken met API's, als om iets nuttigs neer te zetten binnen een organisatie. Het idee dat het resultaat effectief gebruikt kan worden binnen VLAIO en kan uitgroeien tot een volwaardige ondersteuning voor vercshillende stakeholders, hield mij gemotiveerd doorheen het hele traject. Ook het samenwerken met mensen die dagelijks met deze data en rapporten werken, was zeer interessant.\\

Graag wil ik ook een aantal mensen expliciet bedanken. In de eerste plaats gaat mijn dank uit naar mijn co-promotor bij VLAIO, Karel Coenye, voor de inhoudelijke begeleiding en de feedback die ik kreeg. Ook Jade Van den Eynde wil ik bedanken. Als stagebegeleider bracht zij het onderwerp aan en bleef ze betrokken doorheen het traject. Daarnaast wil ik ook nog enkele andere collega’s van het CC-Data team bedanken voor hun hulp en input  tijdens afstemmomenten, alsook de VLAIO-medewerkers die de vragenlijst hebben ingevuld. Hun tijd, voorbeelden en kritische bedenkingen hebben deze bachelorproef rijker en realistischer gemaakt dan ik alleen had kunnen bereiken.\\

Tot slot bedank ik mijn promotor Stijn Lievens voor de academische begeleiding en de vele feedbackmomenten op de structuur en uitwerking van de tekst. Dankzij zijn opvolging doorheen het traject heb ik deze bachelorproef kunnen uitwerken tot iets waar ik trots op ben, ook al was het traject niet altijd even eenvoudig.